% !TeX root = ../main.tex

\chapter{初步結果與預期貢獻}

\section{初步實驗結果}

本研究已初步完成硬體平台搭建與軟體管線（Pipeline）整合，並針對影像品質與自動對焦效能進行了驗證。

\subsection{影像品質與色彩校正驗證}
為驗證影像處理管線的必要性與有效性，我們首先以 NBS 1952 標準解析度測試卡進行明場拍攝。如圖 \ref{fig:resolution_correction} 所示（概念說明），未經校正的原始影像因 Raspberry Pi Camera 與顯微物鏡的主光線角度（CRA）差異，呈現出嚴重的周圍暗角以及拜爾濾波器（Bayer Filter）造成的整體綠色色偏，同時感測器表面的灰塵也清晰可見。在導入平場校正（Lens Shading Table, LST）後，系統成功消除了暗角與色彩串擾，獲得均勻的中性背景，並透過增益圖（Gain map）數位移除了灰塵瑕疵，確保了高頻條紋細節的清晰度。

在螢光模式的驗證中，我們採用了濃度 $1.0 \times 10^6$ beads/mL 的 15 \textmu m 黃綠色螢光微球（FluoSpheres）作為樣本。在未進行色彩校正矩陣（Color Correction Matrix, CCM）轉換前，原始影像中的螢光微球顯得蒼白且飽和度低（Whitish），這是由於消費級感測器的寬頻譜響應與通道間的串擾所致。經過全域 CCM 與空間變異張量解混（Tensor Unmixing）校正後，不僅恢復了螢光微球真實、飽滿的綠色，更達成了優異的空間色彩均一性，徹底消除了視野中心與邊緣間的色彩偏移（Color Shift）。這些結果證明，本研究建立的運算校正管線成功消除了消費級硬體的光學缺陷，使其具備輸出研究級定量影像的能力。

\subsection{自動對焦演算法效能評估}
我們利用 Z 軸堆疊（Z-stacks）的影像序列評估了動態自動對焦演算法的精準度。在明視野環境下，實驗數據顯示全域變異數（Global Variance）的曲線寬廣平滑，非常適合作為粗調範圍的導航；而 Tenengrad/Brenner 指標則呈現出狹窄且尖銳的單一峰值，適合精確鎖定焦平面。

在充滿背景雜訊的螢光模式中，傳統指標往往會產生大量假峰值。實驗結果證明，本研究提出的「形態學 + Top 0.1\% 梯度」指標能有效壓抑雜訊基線，產生清晰且單一的對焦峰。綜合上述演算法，搭配樂高龍門架構的高剛性優勢，本系統在 937.5 \textmu m 的長行程搜尋範圍內，能在 14 秒內穩定收斂並鎖定焦點。此過程並未出現純 3D 列印機構常見的震盪或背隙（Backlash）漂移現象，證實了本系統在實現大範圍、全自動高通量影像掃描上的可行性與可靠度。

\section{預期貢獻}
本研究預期達成以下三項具體貢獻：
\begin{enumerate}
    \item \textbf{開創高彈性之開源硬體平台}：首創將樂高光學桌概念應用於倒立顯微系統，解決傳統 3D 列印顯微鏡架構僵化與擴充困難的問題，提供生物與農業研究人員一個零門檻、高度可重構的光學開發環境。
    \item \textbf{消費級感測器之研究級校正}：建立一套開源之 CRA 校正與張量色彩解混管線，打破高昂科學相機的壟斷，為低成本量化螢光分析奠定基礎。
    \item \textbf{結合前瞻之計算光學技術}：成功整合散斑照明傅立葉疊層技術與高強健性自動對焦，在極低的硬體成本下，實現超越物理繞射極限的高解析度螢光成像，極具潛力應用於微流體監測、水質分析與即時細胞培養觀察等多元領域。
\end{enumerate}